% ===== §8 The Orthogonality of Refinement and Justice =====
\section{Refinement Is Not Direction: The Orthogonality Guard}
\label{p22:sec:orthogonality}

\noindent
The two preceding sections are dangerous if left alone. The genealogy showed the aesthetic and the ethical to be coupled across the tradition; the mechanism showed them to be one perception under two descriptions, the sensibility a single moral-aesthetic organ. Taken without what follows, this invites a conclusion the whole paper must refuse, that to refine the sensibility is to become good, that the cultivated perceiver is by his cultivation the just one, that aesthetic education guarantees ethical formation. The conclusion is false, and its falsity is not a minor qualification but a load-bearing part of the account, because a philosophy of aesthetic education that did not face it would be the aestheticism its critics have always, and often rightly, feared, the doctrine that the fine soul is the good soul and that beauty and virtue rise together by nature. They do not always rise together. A sensibility can be exquisitely refined and exquisitely unjust, and any account of aesthetic education that cannot say how owes its readers an explanation it does not have. This section provides it, and the provision is the same orthogonality the series established for cycles in general, brought to bear on the sensibility itself.

\subsection{The counterexample that will not go away}
\label{p22:sub:orth-counterexample}

The objection has a name in the history of the twentieth century, and it should be stated at full strength rather than softened. There were regimes, and there have been smaller cruelties without number, in which a genuine and highly cultivated aesthetic sensibility coexisted with, and even served, atrocity: the commandant who wept at the cello and returned to his work, the refined taste that ordered beautiful things and unbeautiful deaths with the same discriminating hand. This is not the philistinism of the brute who feels nothing; it is the more disturbing case of a real fineness of perception that does not issue in justice and may even lend its fineness to injustice, dressing domination in a beauty that makes it more seductive and not less. The case will not go away because it is real and because it is precisely the case the coupling account seems to rule out. If the sensibility that perceives the appropriate is one with the sensibility that perceives the good, how is the exquisite torturer possible, who perceives the appropriate with such delicacy and the good not at all? An account that cannot answer has been refuted by the cello.

The case is sharper still at the scale this paper works in, the scale of the intimate field, because there it wears no uniform and commits no atrocity, and so is harder to recognise and more common.

\begin{casebox}[The beautiful house built on one person's disappearance]
Consider a shared life of real and perceptible beauty. The mornings are finely made, the rhythms tended, the placement of things exact, the contingencies caught and shared; the field, by every dynamical measure, accumulates, and the perception that builds it is genuinely refined, not a performance of refinement but a true fineness of attention to the state of the field and the appropriate gesture within it. And one of the two is the fuel. The beauty is sustained by her continual small self-erasures, her preferences quietly dropped, her perceptions deferred, her generativity spent entirely on the field's accumulation and never on her own, so that the spiral the other perceives so finely is driven by the steady burning of her. The other's perception is not false; he sees the field's accumulating beauty truly, and he is not a brute who feels nothing but a man of real sensibility who has simply never been made to perceive that the beauty he attends to so finely is paid for by her disappearance. His fineness is real and bears on the field's accumulation; his blindness is real and bears on who fuels it; and nothing in the first supplies the second. This is the exquisite torturer brought home, and he is not rare. He is the refined partner whose refinement has never once turned to ask who pays for the beauty he perceives so well.
\end{casebox}

This intimate form of the case is the one the paper most needs to meet, because it is the failure its own account of field-building could most easily produce: a cultivation that raised the perception of the field's accumulation to great fineness while leaving wholly untouched the question of who, within the accumulating field, is being consumed to drive it. An aesthetic education that did only what it says it does, refine the perception of the generative direction, would produce exactly this man, fine and unseeing, and it is to forestall him that the orthogonality must be established and built into the cultivation as a separate and standing demand.

\subsection{Two features of the field, not one}
\label{p22:sub:orth-two-features}

The answer requires distinguishing two features of a relational field that the coupling account, stated baldly, runs together, and the distinction is the one the series has drawn before between the dynamical and the just. The first feature is the field's generative direction, whether it accumulates holonomy or dissipates it, whether the path returns to its root raised or closes on its origin flat. This is what the sensibility, refined, perceives finely, and it is what the previous section called the appropriate-and-good seen as one. The second feature is altogether different: it is the question of who, within the field, is rendered mere fuel for its accumulation, whether the holonomy that the field gains is gained by treating one of its members as material to be consumed, a means whose own generativity is extracted to feed another's. A field can accumulate holonomy beautifully, the path spiralling upward, the perception of its accumulation exquisitely fine, while one of its members is the fuel being burned to drive the spiral. The fineness of the perception bears on the first feature, the accumulating character of the field. It is silent on the second, the justice of who pays for the accumulation.

\claim{Refinement of the sensibility and justice of direction are orthogonal. The sensibility, refined, perceives finely the generative direction of the field, whether it accumulates holonomy or dissipates it; this is the aesthetic-and-ethical perception the coupling section established. It does not thereby perceive whether the accumulation is achieved by rendering a member of the field mere fuel, for that is a distinct feature, the justice of the field, which concerns who pays for the accumulation rather than whether accumulation occurs. A field may accumulate beautifully upon the consumption of one of its members, and a sensibility may perceive that accumulation with great fineness while remaining blind, or indifferent, to the consumption. Fineness of perception is therefore no guarantee of justice, and the cultivated sensibility is not by its cultivation the just one.}

This is the orthogonality the series proved for cycles, now located in the perceiver. There it was shown that the dynamical criterion, whether a cycle is generative, and the justice criterion, whether someone is rendered mere fuel, are independent, so that a cycle can be generative and unjust at once. Here the same independence is shown to hold for the sensibility that perceives the cycle: a perception refined upon the generative direction is not thereby a perception of the justice of the field, because the generative direction and the justice are two features, and to be fine upon the one is not to be even awake to the other. The exquisite torturer is possible because his fineness is real and pertains to the first feature, the accumulating character of his field, while his blindness pertains to the second, the fuel his accumulation consumes, and nothing in the refinement of the first supplies the second. The cello is not a fraud. It is a fine perception of a real beauty in a field whose justice the perceiver has never been made to see.

\subsection{The survival of the coupling under the orthogonality}
\label{p22:sub:orth-coupling-survives}

It might seem that the orthogonality destroys the coupling the previous two sections built, that if fineness does not deliver justice then the aesthetic and the ethical have come apart after all. They have not, and seeing why is essential. The coupling was always a coupling of one perception to one register of value, the generative direction perceived as appropriate and as good; it was never the claim that this one perception exhausts the ethical. Justice, the question of who is rendered fuel, is a further ethical matter, not captured by the generative direction and not perceived by the sensibility refined upon it, and the coupling section's good was always the generativity of the field and not the whole of the ethical. So the orthogonality does not break the coupling; it bounds it. The aesthetic and the ethical are one at the level of the generative direction, which the sensibility perceives under both descriptions; they are two at the level of justice, which the sensibility does not perceive at all. The coupling holds within its domain and is silent outside it, and the orthogonality marks exactly where its domain ends, at the line between whether the field accumulates and who pays for the accumulation.

This is why aesthetic education, as this paper conceives it, requires a second cultivation that it does not itself supply, and why the paper has never claimed that the refinement of the sensibility is the whole of moral formation. The refinement makes the perceiver fine upon the generative direction, able to build fields that accumulate and to perceive the accumulation at fine grain; it does not make the perceiver attend to who, in the accumulating field, is being consumed. That attending is a separate discipline, the cultivation of a regard for whether the members of one's field are treated as ends generating their own value or as fuel feeding another's, and the field-builder who has only the first cultivation is the dangerous figure the counterexample named, fine and unjust, the more dangerous for being fine. The paper's account of aesthetic education is therefore deliberately incomplete as an account of moral formation, and the incompleteness is honest rather than evasive: it says exactly what aesthetic education does, raise the resolution upon the generative direction, and exactly what it does not do, supply the regard for justice, and it refuses to let the first masquerade as the second.

\subsection{The guard carried into the practice}
\label{p22:sub:orth-into-practice}

The orthogonality is not only a theoretical safeguard but a practical instruction, and it bears directly on the four-step cultivation. At each step the danger has a specific face. A perception refined upon contingency can be a perception that sees, with exquisite delicacy, the accumulating beauty of a field built upon another's consumption. A reception can welcome a beauty whose production rests on rendering the other fuel. A sharing can be a joint attention between two who are accumulating beautifully upon a third who is being burned. A reproduction can thicken, cycle upon cycle, a good cycle that is unjust at its root. The four-step model, left to its dynamical criterion alone, cultivates the generative direction and is silent on the justice, exactly as the sensibility is, and so the model must carry the orthogonality guard within it as a standing question, asked at every step alongside the question of accumulation: not only whether this field is gaining, but who, if anyone, is being made to pay for the gain. This question is not answered by any refinement of perception, because it is orthogonal to refinement; it must be asked deliberately, as a discipline distinct from and additional to the cultivation of the sensibility, and an aesthetic education that did not install it would produce the fine and the unjust with great efficiency. The guard is the paper's insistence that the cultivation of beauty, even relational beauty, even the beauty of accumulating fields, is necessary and not sufficient for the good, and that the gap between the necessary and the sufficient is exactly the place where justice does its separate, unglamorous, indispensable work.

\subsection{The coexistence of the two cultivations}
\label{p22:sub:orth-coexist}

If the cultivation of the sensibility and the cultivation of the regard for justice are orthogonal, a question arises about how they are to be held together in a single education, since orthogonality might seem to make them two unrelated curricula awkwardly juxtaposed. The question is fair and the answer shows the orthogonality to be a relation between distinct things rather than a separation of unrelated ones. The two cultivations are orthogonal in that neither delivers the other, the refinement of perception not yielding the regard for justice and the regard for justice not yielding the refinement of perception; but they are not unrelated, because they bear on the same field and are exercised in the same acts, the one asking whether the field gains and the other asking who pays for the gain, both questions about the one field the builder is building. They are therefore not two curricula but two questions asked of one practice, distinct in that the answer to one does not give the answer to the other, joined in that both are asked, at every step, of the same field-building act. The education that installs both is not a juxtaposition of unrelated trainings but a single training in which two distinct questions are made habitual, the question of accumulation and the question of justice, asked together of each act though neither is reducible to the other.

This coexistence has a definite structure, and naming it completes the account. The question of accumulation is the question the sensibility answers, by perceiving whether the field gains, and its cultivation is the raising of the perception's resolution. The question of justice is not answered by the sensibility at all, by the orthogonality, and so its cultivation cannot be the raising of the same perception's resolution; it must be the installation of a distinct habit, the habit of asking, of each accumulating field, who within it is rendered fuel, a habit that does not perceive the answer in the way the sensibility perceives accumulation but raises the question that the perception of accumulation would otherwise never raise. The regard for justice is thus a cultivated interruption of the sensibility's absorption in accumulation, a habit of stepping back from the perception of the gaining field to ask the question the gaining field does not pose, and the education that installs it is the cultivation of this stepping-back alongside the cultivation of the perception it interrupts. The two coexist as a perception and its disciplined interruption, the fine seeing of the field's gain and the habitual question that the fine seeing does not contain, and an aesthetic education worthy of the name cultivates both, the resolution that perceives accumulation and the regard that asks, of every accumulation, whether it is just.
